\begin{pastexam}[label={px:ncu-112-dsa-02}]{中央資工所 112 學年度 複選題第 2 題}
（複選；每題 5 分，答錯 1 題倒扣 1 分）

Apply Quick sort on a given list [9, 7, 24, 6, 11, 4, 2, 19]. Which of the following two are the
resulting sequences of the first phase, when the first element or second element is chosen as pivot?

(A)~4, 7, 2, 6, 9, 11, 24, 19\\
(B)~7, 6, 24, 11, 9, 4, 2, 19\\
(C)~2, 6, 4, 7, 11, 24, 9, 19\\
(D)~7, 6, 4, 2, 9, 24, 11, 19
\end{pastexam}

\begin{pastsoln}
本書採用的答案是 (A)、(C)。

題目沒有規定分割（partition）的寫法，而第一趟的結果正是由分割的寫法決定，所以這題有解讀的空間。 下面
先用「分割完成」的條件篩掉不可能的選項，再說明 (A)、(C)、(D) 各是哪一種寫法的結果。 位置從 $1$ 數起。

\textbf{必要條件。} 不管哪一種寫法，第一趟結束時基準 $x$ 已經在最後的位置，左邊都 $< x$、右邊都 $> x$
（資結書第~14.5~節）。 基準只可能是 $9$（第一個）或 $7$（第二個）。
\begin{itemize}
\item (A)：$9$ 在第 $5$ 格，左邊 $4, 7, 2, 6$ 都 $< 9$，右邊 $11, 24, 19$ 都 $> 9$。 可能。
\item (B)：$9$ 在第 $5$ 格，左邊有 $24 > 9$；$7$ 在第 $1$ 格，右邊有 $6 < 7$。 兩個基準都不成立，
\textbf{任何}寫法都不會分割出這個結果，錯。
\item (C)：$7$ 在第 $4$ 格，左邊 $2, 6, 4$ 都 $< 7$，右邊 $11, 24, 9, 19$ 都 $> 7$。 可能（基準 $7$）。
\item (D)：$9$ 在第 $5$ 格，左邊 $7, 6, 4, 2$、右邊 $24, 11, 19$。 可能（基準 $9$）。
\end{itemize}

\textbf{基準取第一個（$9$）：台灣課本常用的兩指標分割。} 這是資結書第~14.5~節作答提醒的寫法：$i$ 從
第 $2$ 格往右找 $\ge x$ 的數，$j$ 從最後往左找 $\le x$ 的數，還沒交錯就交換；交錯之後把基準和 $A[j]$
交換。
\begin{center}
\small
\begin{tabular}{lll}
\toprule
步驟 & 動作 & 陣列 \\
\midrule
開始 & $x = 9$ & $9, 7, 24, 6, 11, 4, 2, 19$ \\
1 & $i$ 停在 $24$（第 $3$ 格），$j$ 跳過 $19$ 停在 $2$（第 $7$ 格），交換 & $9, 7, 2, 6, 11, 4, 24, 19$ \\
2 & $i$ 停在 $11$（第 $5$ 格），$j$ 停在 $4$（第 $6$ 格），交換 & $9, 7, 2, 6, 4, 11, 24, 19$ \\
3 & $i$ 停在 $11$（第 $6$ 格），$j$ 停在 $4$（第 $5$ 格），已交錯 & \\
4 & 基準和 $A[5] = 4$ 交換 & $4, 7, 2, 6, 9, 11, 24, 19$ \\
\bottomrule
\end{tabular}
\end{center}
得到 (A)。 同一種寫法若改取 $7$ 當基準（先和第一個交換成 $7, 9, 24, 6, 11, 4, 2, 19$），結果是
$6, 2, 4, 7, 11, 24, 9, 19$，不在選項裡。

\textbf{基準取第二個（$7$）：Lomuto 式的單向掃描。} 先把基準和第一個對調（下表第一列）。
$i$ 指向「$< x$ 區」的最後一格（一開始是基準那格），$j$ 從第 $2$ 格往右掃，$A[j] < x$ 就把 $i$ 加一、
交換 $A[i]$ 和 $A[j]$（資結書虛擬碼~14.8 的想法，只是基準放在最前面）。
\begin{center}
\small
\begin{tabular}{lll}
\toprule
$j$ & 動作 & 陣列 \\
\midrule
開始 & $x = 7$，和 $9$ 對調；$i = 1$ & $7, 9, 24, 6, 11, 4, 2, 19$ \\
2、3 & $9$、$24$ 都 $> 7$，不動 & \\
4 & $6 < 7$：$i = 2$，交換第 $2$、$4$ 格 & $7, 6, 24, 9, 11, 4, 2, 19$ \\
5 & $11 > 7$，不動 & \\
6 & $4 < 7$：$i = 3$，交換第 $3$、$6$ 格 & $7, 6, 4, 9, 11, 24, 2, 19$ \\
7 & $2 < 7$：$i = 4$，交換第 $4$、$7$ 格 & $7, 6, 4, 2, 11, 24, 9, 19$ \\
8 & $19 > 7$，不動 & \\
最後 & 基準和 $A[i] = A[4]$ 交換 & $2, 6, 4, 7, 11, 24, 9, 19$ \\
\bottomrule
\end{tabular}
\end{center}
得到 (C)。 同一種寫法取 $9$ 當基準，結果是 $2, 7, 6, 4, 9, 24, 11, 19$，也不在選項裡。

\textbf{(D) 的來歷。} 把小於基準的數依原順序放左邊、大於的依原順序放右邊（不是原地分割，要另開陣列），
基準 $9$ 得到 $7, 6, 4, 2, 9, 24, 11, 19$，就是 (D)。

\textbf{結論。} 每一種寫法只產生一個選項，所以「兩個」答案是「基準 $9$ 一個、基準 $7$ 一個」。 基準 $7$
只有 (C) 可能；基準 $9$ 的 (A)、(D) 之中，(A) 是課本的原地分割（題目說 Quick sort，通常指原地的寫法），
(D) 要另開陣列。 所以採用 (A)、(C)。 若以不原地的寫法為準，答案會變成 (C)、(D)。
\end{pastsoln}
